%!TEX root = ../../main.tex
%% 1.2 선행 연구에서 학위논문으로의 개선 논리 (I1–I4) — 근거: MASTER_THESIS_RESEARCH_PLAN §3, COMMON_RESEARCH_SPINE §2
\section{선행 연구에서 학위논문으로: 네 단계의 개선}
\label{sec:intro-improvement}

선행 연구 \cite{baek2026killconditioned}의 각 구성 요소는 표 \ref{tab:cog-to-thesis}의 네 단계 I1--I4로 확장된다. 이 표는 선행 연구가 이미 한 일(고정 계수 교환가치의 한계 명시, 같은 테이블 입력의 MLP 비교, 시간순 패치 분할)을 새 기여처럼 주장하지 않기 위한 장치이기도 하다.

\begin{table}[htbp]
  \centering
  \caption{선행 연구의 출발점과 학위논문의 확장}
  \label{tab:cog-to-thesis}
  \small
  \begin{tabularx}{\linewidth}{@{}l>{\raggedright\arraybackslash}p{4.3cm}>{\raggedright\arraybackslash}p{4.3cm}X@{}}
    \toprule
    단계 & 선행 연구 (CoG 2026) & 학위논문의 확장 & 남기는 근거 \\
    \midrule
    I1 교전 단위 & 손으로 정한 시간·공간 상수의 킬 군집 검출 & 자료에서 추정한 경계 상수, 게임 규칙에 고정한 게이트 상수, 상수별 출처 분류, 규모 계급 & 상수표, 사례 변화·중첩·제외 분석, 정의 민감도 \\
    I2 결과 가치 & 고정 계수의 교환가치 (교전 국소 라벨) & 경기 승패로 학습해 동결한 평가기 $\Vhat$의 구간 변화 방향 $\SVI$ & $\Vhat$ 품질, 결과 시점 규칙, 분해·안정성·대응 분석 \\
    I3 사전 예측 & 표현·학습기 결합 비교, 정합 입력 MLP & 동일 352개 입력의 고정 사양 예측기 $q$와 승률·시간 기준선 & 코호트별 주 대비, 선정·보정 기록, 학습기 비교의 한계 명시 \\
    I4 의미·범위 검증 & 초기 골드·시간 층화, 인접 패치 평가 & 균형 구간 $\Bforty$, 확률 품질 분해, 작은 변화 제외, 외부 지역·패치 전이 & 성능의 출처·조건·범위에 관한 제한된 결론 \\
    \bottomrule
  \end{tabularx}
\end{table}

그림 \ref{fig:pipeline}은 네 단계가 실제 파이프라인에서 어떻게 이어지는지 보인다.

%% figure 객체: fig:pipeline — 측정·예측·검증 사슬 (TikZ; 값은 본문 매크로). 저널 fig1_pipeline.svg 와 같은 기전.
\begin{figure}[htbp]
  \centering
  \begin{tikzpicture}[
    node distance=5mm and 6mm,
    box/.style={draw, rounded corners=2pt, align=center, font=\small, inner sep=4pt, minimum height=9mm},
    frozen/.style={box, thick},
    arr/.style={-{Latex[length=2mm]}, semithick}
  ]
    \node[box] (tele) {Match-V5 timeline\\프레임 60 s · 사건 ms};
    \node[box, right=of tele] (det) {검출기 v3.3\\$G$, $D$ (자료), $R$, $B$ (규칙), $M$};
    \node[box, right=of det] (coh) {규모 계급 $\nmin$\\T ($\ge 4$) · S (2--3) · pick 제외};
    \node[frozen, below=of det] (V) {동결 평가기 $\Vhat$ (fit85 MLP)\\TRAIN 행은 fold 평가기 $\Vfold{k}$};
    \node[box, left=of V] (pre) {$\Spre$ at $\tcut = K - 15$ s\\$\ppre = \Vhat(\Spre)$};
    \node[box, right=of V] (post) {$S(\eh)$, $\eh = \min(L + 90\,\mathrm{s}, \ldots)$\\$\Ysvi = \ind[\dV > 0]$};
    \node[box, below=of pre] (q) {$q(\xpre)$: 로지스틱, \dimInputQ{} 입력\\코호트별 $q$, $\qS$};
    \node[box, below=of post] (pt) {$\PTflex$: $(\ppre, t)$ 스플라인\\Q\_SELECT에서 선정 후 동결};
    \node[box, below=of V, yshift=-4mm] (contrast) {주 대비 (코호트당 1개)\\$\dBrier(q - \PTflex)$, 경기 군집 부트스트랩};
    \node[box, below=of contrast, text width=0.86\linewidth] (verify) {검증 (추정 대상을 바꾸지 않음): $V\to W$ 품질 · 방향/평균/규모 분리 · 비교전 대조 · 물질·오브젝트 대응 · 지평 · $\Bforty$ · $\lambda s_Q$ · CORP · 외부 점수 전용};
    \draw[arr] (tele) -- (det);
    \draw[arr] (det) -- (coh);
    \draw[arr] (det) -- (V);
    \draw[arr] (V) -- (pre);
    \draw[arr] (V) -- (post);
    \draw[arr] (pre) -- (q);
    \draw[arr] (pre) -- (pt);
    \draw[arr] (post) -- (contrast);
    \draw[arr] (q) -- (contrast);
    \draw[arr] (pt) -- (contrast);
    \draw[arr] (contrast) -- (verify);
  \end{tikzpicture}
  \caption{측정·예측·검증 사슬. 검출기가 공개 timeline에서 교전을 찾고, 동결 평가기 $\Vhat$가 첫 킬 15초 전의 상태와 마지막 킬에 앵커된 결과 시점의 상태를 점수화하며, 그 차이의 부호가 라벨이다. 교전 전 상태의 로지스틱 모델을 승률·시간 스플라인 기준선과 코호트별로 대비한다. TRAIN 행은 fold 평가기가, 나머지는 동결 번들이 라벨링한다. 검증은 라벨이 재는 것을 한정할 뿐 추정 대상을 바꾸지 않는다.}
  \label{fig:pipeline}
\end{figure}


I1의 사례와 시간 경계는 입력 $X$와 결과 $Y$ 모두에 연결된다. I2는 I1이 정한 사례와 구간에 $\Vhat$를 적용한다. I4는 $q$만이 아니라 I1과 I2도 평가한다. 승률·시간 기준선(PT)은 역할상 I4의 비교 기준이지만 실행상으로는 I3의 학습·보정·동결 절차에 함께 들어간다. ``검증''이 뒤에 놓인 것은 결과를 본 뒤 검증 기준을 고른다는 뜻이 아니다. 대상, 지표, 하위집단, 동결 규칙은 해당 실행의 평가 전에 정의되었으며, 과거의 시험 패치 노출은 제 \ref{sec:intro-epistemic} 절에서 그대로 밝힌다.
